\usepackage{amsmath}
\usepackage{amssymb}
\usepackage{mathtools}
\usepackage{bm}

% per inserire uno spazio "fantasma" nella definizione di un'abbreviazione
\usepackage{xspace}

% per inserire un DOI senza problemi coi caratteri "strani" ivi presenti
\usepackage{doi}
\renewcommand{\doitext}{DOI }% originally was "doi:"

% per inserire correttamente le unità di misura SI (incluse quelle binarie)
\usepackage[binary-units=true]{siunitx-v2}
% se si desidera usare / invece che la potenza -1 per indicare "al secondo"
\sisetup{per-mode=symbol}

% per inserire codice di programmazione complesso
\usepackage{listings}% per inserire codice di programmazione complesso
\lstset{
basicstyle=\ttfamily,
columns=fullflexible,
xleftmargin=3ex,
breaklines,
breakatwhitespace,
escapechar=`
}

% modify some page parameters
\setlength{\parskip}{\medskipamount}
\advance\voffset -5mm
\advance\textheight 30mm

% riga orizzontale
\newcommand{\HRule}{\rule{\linewidth}{0.2mm}}
% esempio di creazione di semplici abbreviazioni
\newcommand{\ltx}{\LaTeX\xspace}
\newcommand{\txw}{TeXworks\xspace}
\newcommand{\mik}{MikTex\xspace}
\newcommand{\html}{HTML\xspace}
\newcommand{\xhtml}{XHTML\xspace}

% esempio di creazione di un'abbreviazione con un parametro (il cui uso è indicato da #1)
\newcommand{\cmd}[1]{\texttt{#1}\xspace}
% per citare un RFC, es. \rfc{822}
\newcommand{\rfc}[1]{RFC-#1\xspace}
% per citare un file (es. \file{autoexec.bat}) o una URI fittizia (es. \file{http://www.lioy.it/})
% per le URI vere usare \url o \href
\newcommand{\file}[1]{\texttt{#1}\xspace}
% per inserire codice di esempio in-line
\newcommand{\code}[1]{\lstinline|#1|}
% importante per i pathname Windows perché non si può usare \ essendo un carattere riservato di Latex
\newcommand{\bs}{\textbackslash}
% definizione di un termine: formattazione ed inserimento nell'indice
\newcommand{\tdef}[1]{\textit{#1}\index{#1}}
% meta-termine, usato tipicamente nelle definizioni dei tag
\newcommand{\meta}[1]{\textit{#1}}

% Todo command
\usepackage{xcolor}
\newcommand{\todo}[1]{\textcolor{red}{\textbf{TODO: #1}}}

% TikZ and PGFPlots definitions for foundations
\usepackage{tikz}
\usepackage{pgfplots}
\pgfplotsset{compat=1.18}
\usetikzlibrary{calc, positioning, arrows.meta, shapes.geometric, decorations.pathmorphing, shapes.misc, fit}

\tikzset{
    foundation/.style={
        font=\small\sffamily,
        >=stealth,
        block/.style={draw, rectangle, rounded corners, minimum width=2.5cm, minimum height=1cm, text centered, fill=blue!5, draw=blue!80!black, thick},
        implicit_block/.style={block, fill=red!5, draw=red!80!black, densely dashed},
        arrow/.style={thick, ->},
        axis/.style={thick, ->},
        exact_curve/.style={thick, black},
        approx_curve/.style={thick, red, dashed},
        symp_curve/.style={thick, blue},
        vector/.style={thick, -stealth, blue},
        force/.style={thick, -stealth, red},
        velocity/.style={thick, -stealth, green!60!black},
        particle/.style={circle, fill=black, inner sep=1.5pt}
    }
}